% ============================================================
%  Harald Høffding, Human Thought: Its Forms and Its Tasks
%  (Den menneskelige Tanke, dens Former og dens Opgaver)
%  København og Kristiania: Gyldendalske Boghandel, Nordisk Forlag, 1910.
%  ENGLISH TRANSLATION
%  Source of truth: transcription.tex (Danish), COMPLETE (392 pp.).
%  Page-offset (scan only): PDF = printed + 15.
%
%  HOW TO USE THIS FILE
%  - This is a SKELETON. Every heading is in place; the prose is not.
%  - `grep -n "text to be added" translation.tex` lists the markers in order.
%  - For each marker: open transcription.tex, find the SAME heading/label,
%    read from that heading to the next, translate the span, replace the marker.
%  - Conventions: ../../../TRANSLATION-PLAYBOOK.md + ./TRANSLATION-RESUME-NOTES.md
%  - Heading glosses below are PROVISIONAL — verify against the Danish on first
%    use (see TRANSLATION-RESUME-NOTES.md "Heading glosses" for flagged choices).
%  - Markers under a heading that has NO intro text before its first sub-heading
%    (B., C. in Ch. I; A., B. in Ch. IV; a. Det etiske) cover an empty span —
%    just delete that marker.
%  - The user commits and pushes. You never do.
% ============================================================
\documentclass[12pt,a4paper]{book}
\usepackage[utf8]{inputenc}
\usepackage[T1]{fontenc}
\usepackage{libertinus}
\usepackage{libertinust1math}
\usepackage{textalpha}   % Greek glyphs carried verbatim from the Danish
\usepackage[english]{babel}
\usepackage[protrusion=true,expansion=true]{microtype}
\usepackage{setspace}
\usepackage[margin=1.2in]{geometry}
\usepackage{enumitem}
\usepackage{fancyhdr}
\setlength{\headheight}{28pt}
\addtolength{\topmargin}{-13.5pt}
\usepackage{hyperref}

\hypersetup{
  pdftitle={Human Thought: Its Forms and Its Tasks},
  pdfauthor={Harald Høffding},
  hidelinks
}

\pagestyle{fancy}
\fancyhf{}
\fancyhead[LE,RO]{\thepage}
\fancyhead[RE]{\textit{Human Thought}}
\fancyhead[LO]{\textit{\leftmark}}

\setstretch{1.3}

\title{\textbf{Human Thought}\\[6pt]
  \large Its Forms and Its Tasks}
\author{Harald Høffding\\[6pt]
  \normalsize translated from \textit{Den menneskelige Tanke} (1910)}
\date{København og Kristiania:\\ Gyldendalske Boghandel, Nordisk Forlag, 1910}

\begin{document}
\maketitle
\thispagestyle{empty}
\newpage

\tableofcontents
\newpage

%% ============================================================
%%  FOREWORD (Forord, August 1910)
%% ============================================================
\chapter*{Foreword}
\addcontentsline{toc}{chapter}{Foreword}
\label{ch:forord}

\noindent The life of thought, which it is the task of my book to
describe, is intimately interwoven with the other sides of the life of
the spirit. I have sought to show that thought cannot serve the life of
the spirit as a whole unless it first and foremost becomes master in
its own domain and unfolds itself according to its own laws as an
independent side of the life of the spirit. That this unfolding, when it
is rightly carried through, will bring with it great and far-reaching
changes in personal life as a whole, is something of which I become more
and more convinced. There are here oppositions and problems that will
take on ever sharper forms, and I feel no admiration for those who here
sense no sting, although they have sufficient opportunity to come to
know the life of thought.

Just as I have on the whole stressed the personal factor in all
thinking, the more this approaches the border-regions, so too I have not
concealed that my own stance at thought's final station has a personal
character, and I have not been able to refrain from stating how, in my
experience, one feels at this station.

I still have much to learn. My work will need filling out and
correction on many points. But beyond the framework I have given here, I
cannot go. And it affords me a good place to work. May there be one or
another who could also find work here, even if this work should lead him
to alter both the framework and its content.

\bigskip
\noindent August 1910.\hfill \emph{Harald Høffding.}

%% ============================================================
%%  I. THE PSYCHOLOGY OF THOUGHT (FUNCTIONS)
%% ============================================================
\chapter*{I. The Psychology of Thought (Functions)}
\addcontentsline{toc}{chapter}{I. The Psychology of Thought (Functions)}
\markboth{I. The Psychology of Thought}{I. The Psychology of Thought}
\label{ch:psykologi}

\begin{center}---\end{center}

\section*{A. Psychic Energy}
\addcontentsline{toc}{section}{A. Psychic Energy}
\label{sec:psykisk-energi}

\noindent 1. The task I have set myself in this book is to investigate
human thought---a task that all philosophers have set themselves,
whether they have been conscious of it or not. For all questions that
can be raised concerning life and existence depend---both with respect
to the way in which they are posed and with respect to the way in which
they are answered---on the conception one has of the nature and
operation of thought. It is therefore also an outline of my philosophy
that I shall give here. That is to say: I shall give an account of the
thoughts to which I have come back again and again, when I have occupied
myself with such questions as: the presuppositions of science, the
significance of the results of science for one's view of life, and the
relation between scientific cognition and other sides of the human life
of the spirit.

2. We know thought best as reflection. The word \emph{Eftertanke}---``after-thought''---I
use not only in its strict sense, ``in one's thinking to \emph{seek
after} something or \emph{pursue} something,'' but also in the sense
that thought \emph{comes after} an immediate experience. Reflection thus
presupposes that we have seen, felt, experienced something. During the
experience we were perhaps wholly absorbed in it; but now we seek to
bring it forth anew, either to dwell once more on what was experienced,
or to go through its various parts or properties or its relations to
other experiences. In mere dwelling, reflection is one with a beholding,
a mere representing. But such dwelling can scarcely persist for long in a
wholly quiescent form. Involuntarily---perhaps precisely because of the
interest that what is experienced holds---a movement of thought will be
released, one that serves to bring out its inner and outer relations.
Thus when the poet involuntarily reaches for similes and images in order
really to let his subject come to expression. Or when perhaps it is
precisely the contrast with other experiences that serves to illuminate
the peculiar character of the single experience. Both likeness and
difference can thus serve to cast a clearer light on what was given in
the original experience, and reflection therefore does not necessarily
stand in opposition to it, but can work in its service and in its
spirit, when it has itself vanished.

By the word experience we think chiefly of states in which we were
essentially receptive. But reflection can also serve the involuntary
life of the spirit, where this appears as urge or striving. In an
instinctive or inspired manner, involuntary striving is often led to its
goal without any reflection being needed, and perhaps better than could
happen with the help of any reflection. When reflection comes to operate
in our action, it will be because the involuntary striving met
resistance. There may then be needed a deliberation upon what it really
is that is being striven for, upon this goal's relation to other
possible or actual endeavors, and upon the ways and means that can lead
to the goal. Here too it is through the bringing-out of likenesses and
differences that reflection renders its service. And here too a
difference can appear between dwelling and movement. It may be that the
representation of the goal is held fast without great oscillations, and
that the relation to ways and means or to other possibilities plays a
subordinate role, as when the helmsman stares at the compass or the
seamark, and his hand involuntarily moves the rudder according to what
he sees. But it may also be that the strongest movement back and forth
among various possibilities takes place and must take place before the
decisive concentration that is called decision sets in.

The arising of reflection is like a kind of awakening, not from sleep or
from unconsciousness, but from the involuntary mental life and its
experiences, from states that may provisionally be characterized as
sensation, intuiting, representing, mood, emotion, urge, and striving.
In its dwelling form, reflection stands closest to the involuntary
mental life. In its striving to bring out likenesses and differences, it
moves further away from the involuntary. Linguistic usage marks this
difference rather aptly by distinguishing between \emph{thinking of}
something and \emph{thinking over} something. It is the latter kind of
reflection, the discursive, mobile kind, that we shall have most use for
in what follows, since it is the one we can best observe, and likewise
the one that plays the decisive role in the tasks of thought. Only for
this kind of reflection do problems exist. In what follows I therefore
understand by thinking this kind of reflection. And yet there is only a
difference of degree between the two kinds. Only for a brief moment is a
dwelling representing, a holding-fast of a memory-image, possible;
precise observation shows that in greater or smaller rhythms we move
toward and away from the object we are thinking of. And every return can
bring with it a recognition, so that the relation of likeness comes to
play a role just as in discursive reflection. Perhaps the representation
was new to us at the first moment; but through repeated dwelling it
becomes familiar. Apart from this too, dwelling reflection by no means
excludes that manifold other representations and experiences can
contribute and perhaps precisely condition the interest that makes
continued holding-fast of a single representation possible. What holds
for discursive reflection will therefore also hold, to a certain degree,
for the dwelling kind.

It emerges from this description that the work of thought consists in
comparison, in distinguishing or finding likenesses. When I think over
something, I seek to determine its relation to something else, and the
more relations of likeness and difference I find, the better I have
thought it through. This holds of course also when it is one and the
same object that I examine in different connections and under different
conditions; I then compare the different ways in which the object
appears in the different cases.

But comparison in turn presupposes a holding-together or combining,
synthesis. I can compare two objects only when I either have them given
to me simultaneously in my sensation, my memory, or my imagination, or
else retain a representation of the one while the other is present to
sensory observation. If the one object constantly vanished entirely when
the other emerged, no comparison could take place. The work of thought
consists in a combining by which likenesses and differences become
determinative for the relations in which the objects are placed to one
another in the connection that results from the work. Synthesis is an
expression of energy, of the capacity to perform work.

The resistance that is overcome through this work is due to the objects
against which reflection is directed and which are given in the
involuntary mental life---in sense-sensation, memory or imagination,
feelings of pleasure and displeasure, urge and striving.

The combining work must become the more difficult, the greater the
multiplicity of objects that are to be brought into relation with one
another. The fullness of content can be so great that the psychic energy
available in the particular case cannot master it. There then threatens
a dissolution of the mental life, a transition to chaos.

Yet it depends not only on the multiplicity, but also on the degree of
difference that the objects present. The more they resemble one another,
so that they can fuse, or so that at least the work done in the case of
the one need not be repeated in the case of the other, the less work is
required. Objects that stand in strong mutual opposition, or perhaps
even conflict, require a great energy if they are all to be taken up into
the conscious connection of the mental life. This holds both for the
inquirer, who must find a point of view from which the conflicting can
stand as one, and for the practically striving person, who must
subordinate the conflicting impulses and interests to the thought of a
great and guiding purpose.

It will, thirdly, depend on how independent a place the individual
objects are to occupy in the connection that is being worked on. A
schematic whole is easier to form than a whole that lets the
peculiarities of the individual objects each come into their own. Here
it is precisely likeness that can create a difficulty, since it easily
leads to fusion and thereby to the abolition of the independence of
individual objects. On the other hand, from this point of view a
pronounced difference or opposition can be a relief.

Furthermore, a greater work is required the further apart the objects
lie from one another in time and space. The simultaneous and the near
join together involuntarily, and the connection perhaps even forces
itself upon us. And if we leave the objects before us in order to
immerse ourselves in the world of memory and imagination, here too the
connection can often form itself with ease; ``with great ease thoughts
take up their abode with one another,'' but when it is a matter of
holding together a present object with distant memories, or with ideals
that have been properly to the fore in the soul only in isolated moments
of enthusiasm, then special demands are made on psychic energy. To have,
in the situation of the moment, the treasures of one's experiences,
one's reflection, and one's ideality fully at one's disposal is perhaps
the greatest testimony to mental energy.

Finally, it plays a role how close and intimate a connection the
objects, despite their difference, are to enter into. There are here
many degrees, from wholly external juxtaposition and grouping to the
inner connection that a scientific concept, an artistically formed
image, or a great decision can express.

Besides these five conditions---which might be said to form an analogy
to what mass is in the case of physical energy---account must also be
taken of the speed with which the synthesis proceeds. The genius and the
man of will can in an instant expend a capital of energy that other
people must distribute over a longer span of time. In the domain of the
life of feeling, this opposition appears as an opposition between
vehemence and inwardness, emotion and disposition\footnote{See on this
my \emph{Psykologi} IV, 7; VI E, 4--5.}, and analogous oppositions hold
for the other sides of the mental life. In research, discovery often
forms an analogy to emotion, the conduct of proof to disposition.

The psychic energy that expresses itself in synthesis we measure by
taking all these circumstances into account. Of exact measurement there
is of course no question here, only of an evaluative estimate. But as
soon as we seek to give grounds for why one work of thought, one
intellectual orientation, one view of life can be called lower or higher
than another, we must take all these circumstances into account, and we
in fact do so, more or less consciously, and more or less precisely.
There is less psychic energy available in the dream-state, where the
changing objects each release their own train of representations, than
in the waking state, where dominant purposive thoughts bring about a
main direction. It is a purely psychological measure that we apply here.
And even though, in forming the concept of psychic energy, we have had
guidance in the analogy with the concept of physical energy, that
concept is nevertheless built upon its own peculiar experiences, in
which physical or physiological observations and investigations need
play no decisive role. For this reason we need not fail to appreciate
the significance of the attempts that have been made, from the
experimental-psychological and physiological side, to measure the
brain-work presupposed in certain elementary mental activities---e.g.
calculating or remembering---through its relation to the muscular work
that can be performed simultaneously with it. One can detect such
thought-work by the diminution it produces in the muscular work
performed at the same time. But this very measurement presupposes that
one can, by purely psychological means, distinguish between more
difficult and easier mental work; otherwise one moves in a circle. The
brain-work corresponding to the mental work described above is, after
all, not known directly, but only through the very mental work to which
it is assumed to correspond.

3. Reflection presupposes objects on which it can exercise its combining
and comparing activity. But now these objects? Does the law of synthesis
hold for them as well? We stand here before a question of a peculiar
difficulty, closely bound up with the whole nature of the mental life.
It is only reflection that acquaints us with the objects. It is the
medium through which they are accessible to us. The involuntary mental
life in its immediate unfolding must not be described and judged
straightway according to the way in which it shows itself to reflection.
It is a life within which we are not permitted to presuppose properties
and determinations that first come into being through the work of
reflection. When, in isolated moments, we are able to immerse ourselves
in this involuntary---but not for that reason unconscious---life in such
a way that we can afterwards retain a memory of it, it stands before us
as a stream, a gliding-along of manifold objects, which we designate as
sensations, memories, moods, endeavors, and so on. It is not a chaos of
independent objects; already in the involuntary mental life there are
connections and wholes as well as differences and oppositions. Every
word we use about what we can retain from the involuntary states bears
witness to this. Most striking, no doubt, is the continuity; this is why
the word ``stream'' offers itself so readily for what we call the
immediately given. But it is not a uniform stream. Differences of
quality make themselves felt; otherwise we could not distinguish between
the kinds of objects called above sensations, memories, and so on. And
there can be eddies in the stream, or windings, so that the direction is
not kept in a straight line. That is to say, strivings of different
tendency, of more or less opposite direction, can make themselves felt.
With this is bound up, again, that inhibitions can set in, while in
other cases, where the different tendencies go in the same direction, a
facilitation or an acceleration can appear.

Through the indefinite and figurative expressions that alone stand at
our disposal for indicating the basis of objects that reflection
presupposes, it shines through that this basis forms no absolute
opposition to what emerges as a consequence of the work of reflection.
There too we find both unity and multiplicity, both continuity and
discontinuity, and we find many degrees and nuances of these relations.
It is therefore not justified, with the older English school
(\emph{Locke}, \emph{Hume}, \emph{Mill}), to conceive the immediately
given in the mental life in an atomistic way, as if it consisted in a
multiplicity of mutually independent elements. Both in the psychic and
in the physical domain, it is only reflection that leads to the
formation of atomic concepts, guided by the urge to make the differences
as distinct as possible, so that the connection of them can become as
clear and perspicuous as possible. The unity to which reflection bears
witness, and which it itself seeks to bring about, thereby becomes
enigmatic. But it is also unjustified, as in our day especially
\emph{William James}\footnote{\emph{The Stream of Thought}. Mind 1884
(reprinted as Chap.\ 9 of \emph{Principles of Psychology}). Later
especially \emph{A pluralistic universe} (1909), Sect.\ 7.} and
\emph{Henri Bergson}\footnote{\emph{Les données immédiates de la
conscience} (1888). Later especially in \emph{Introduction à la
Métaphysique}. (Revue de Métaphysique et de Morale 1903).} have
attempted, to present that involuntary life as standing in the sharpest
opposition to reflection and its work. It then becomes, in the last
analysis, the indescribable and inconceivable, and it would become an
insoluble enigma how there can be any knowledge at all of that
``immediately given'' from which reflection is nonetheless constantly to
draw. James and Bergson represent a significant reaction against
intellectualism's misunderstanding and overvaluation of the significance
of reflection, a reaction that in several respects recalls the stand
taken by \emph{Hamann}, \emph{Herder}, and \emph{Jacobi} against
\emph{Kant}. Not pure reason, but the immediate stream of the life of
the spirit is, they maintain, the highest, and brings truth forth. The
distinctions and definitions of thought are necessary in order that we
may practically find our way in life. And analyzing science has carried
the opposition to the immediately given still further---an opposition
that language itself already favors, by forming independent words for
the individual objects and by taking spatial relations as the basis for
our intuitive images. It is as though a Fall had taken place when
reflection awoke and turned, inquiringly, toward the involuntary life
from which it had itself sprung.

Without a kinship between the involuntary mental life and reflection, it
would be impossible to understand how the latter could arise. But in the
very descriptions that James and Bergson give of what, for them, in
sharp opposition to all elaboration by thought, is the immediately
given---descriptions that are masterpieces of descriptive
psychology---it is nonetheless expressly stated that in this ``given''
there are differences and connections, qualities and continuities. James
emphasizes especially the constant transitions (continuous transitions)
and the difficulty of distinguishing between relation and term (relation
and matter related); Bergson stresses especially the undivided and yet
heterogeneous nature of the inner stream (durée hétérogène et
indistincte; durée dont les moments se pénètrent). But with this an
inner possibility for the arising of reflection has been indicated. The
whole weight must not be placed on the external causes (practical need,
the mechanics of language, and the technique of science). The motive for
the awakening of reflection is given with the opposition between
continuity and qualitative multiplicity, between the connection and the
individual elements. Herein lies the possibility of an inhibition of the
involuntary stream, an inhibition that only discursive reflection can
remedy, through the clarity and definiteness it is able to bring. As a
combining that is accomplished through comparison, reflection can bring
about clarity concerning the relation between continuity and
discontinuity, unity and multiplicity. And it is not anything wholly
foreign that it brings in by its work; for already in the immediately
given there appear, after all, likenesses and differences, connection
and multiplicity---and likewise, what is especially significant, an
involuntary striving to bring about harmony among the different
tendencies. If there were not such an original striving, it would be
unjustified to speak of a stream, of a continuité mouvante, a unité de
direction\footnote{One of my young listeners wrote to me some years ago:
``I can declare myself free of every purpose, as of every interest or
belief in anything outside or inside me. I do indeed feel something
willing, something pressing forward, but it lives, as it were, its own
life, pushes aside everything that gets in its way, and surely never
commits any act that is not to its own advantage.''---Here reflection
has brought no result, only dissolution; but the involuntary life has
maintained itself in its unity and in its definite direction, which was
all the more striking to the young psychologist the less his reflection
could attain to any firm standpoint.}! There shows itself here, then, an
analogy to what, in the case of reflection, we called synthesis---a
tendency to gather and unite. And reflection awakens when there is a
special urge for a fuller carrying-through of this tendency than is
possible within the involuntary mental life. Already within this---not
only after reflection has awakened---crises and catastrophes can arise,
and chaos can threaten. A work can be required even before reflection
can awaken; but the difficulties arise, are worked on, and are resolved
on different terms in the involuntary mental life than where reflection
has taken the helm. Perhaps some of the most significant things in our
mental life take place in the regions below the domain of reflection,
and it is often only the after-effects of these deep-lying processes
that reflection has to deal with. It will always be impossible to
demonstrate that at any point we stand before an immediately given in
the absolute sense, that is, one in which absolutely no psychic work has
been done. What we in each particular case call immediately given
deserves this name only in relation to a certain definite psychic work,
while in comparison with other objects it would perhaps have to be
regarded as the result of a work.

4. By starting from reflection and seeking to become clear about its way
of operating, we arrived at the concept of synthesis, and we then found
that this concept could also be applied to the involuntary mental life,
which forms the constant basis from which reflection springs, in order
thereupon to make that basis itself an object of its activity. We must
make ourselves familiar with the fact that the involuntary mental life,
as it unfolds prior to the awakening of reflection, can be illuminated
and designated only by way of analogy. We apprehend it only in analogy
with clear reflection, and the expressions of language are throughout
drawn from the world of reflection. So far is that involuntary mental
life from being the thing easiest to apprehend, that it takes precisely
great art to get a proper view of it in its difference from the mental
life shaped by reflection. What one calls sound common sense operates,
in its psychology, with concepts drawn from the world of reflection. A
childlike tendency makes itself felt to regard everything in the mental
life as the fruit of clearly conscious thinking, of deliberation and
calculation. One personifies, so to speak, not only outer but also inner
nature, placing a person behind the person. The childlike psychology
explains personal life itself, in its involuntary unfolding, by
supposing that behind it there is a person who consciously thinks
whatever stirs in it---just as in older works on optics one can find
depicted a man inside the head, looking at the image that the light-rays
form on the retina!

Reflection is required in order to discover and apprehend the
involuntary mental life, and yet the peculiarity of the latter consists
in the fact that reflection plays no role in it. This difficulty can be
resolved only by our apprehending and expressing the
involuntary---insofar as any memory of it persists after the awakening
of reflection---in analogy with what we know clearly from the world of
reflection itself. This is what we do when we say that synthesis is the
fundamental form, not only for comparing reflection, but also for the
involuntary mental life, and is thus the most comprehensive
psychological concept.---

But now a new difficulty arises. Synthesis means combining---but there
must then be something that is combined, and this something must be a
multiplicity. The objects toward which reflection directs itself in
order to hold them together and compare them form a coherent
multiplicity; each individual object comes into being for us through an
act of distinguishing. And in a single object we can again distinguish
individual sides, properties, or parts, which we may call elements. The
result seems then to be that we end in an opposition between synthesis
and elements, and that the multiplicity of elements is the constant
presupposition for synthesis being able to operate! From this one might
then draw the conclusion that what is original in the mental life was a
multiplicity of elements, and that only at a later point of development
did the combining activity intervene. This conception was also asserted
in its day by John Locke, and it still often haunts psychology. It
relies on the analogy with the formation of a mechanical whole out of
elements that could be pointed out beforehand each by itself, as when a
heap of sand is formed from grains of sand that previously lay each in
its own place and are only now brought together. The question becomes
whether this analogy is justified.

When reflection awakens, the mental life does not, however, begin with
it; for then it would have no objects. The elements---that is, the
parts, sides, or properties that reflection can distinguish in the
object or objects, and that it can determine through their mutual
likenesses and differences---are, after all, precisely the fruit of a
work of thought. They are drawn out of the connection in which they
involuntarily occurred, and set in a certain opposition to one
another---an opposition without which they could not become objects of
clear consciousness. And the characterization that each individual
element, and thereby each individual object, receives is conditioned by
its relation to other elements or objects. At no point is it justified
to assume isolated elements in the mental life, each absolutely
independent by itself---a kind of absolute psychic atoms. That is to
confuse the effect of reflection with reflection's own basis---at
bottom, again an example of the naive psychology that tends, after all,
to confuse differences that thought first makes with originally given
differences.

And insofar as we succeed in retaining a memory of the state we are in
when involuntariness prevails, this state appears---and here the masters
of descriptive psychology are undoubtedly right---not as a chaos of
independent parts, but as a stream, a pulsation, yet not without
differences, but such that there is a connection even between the lowest
troughs and the highest crests of the waves. We find, as far as can be
judged, unity and multiplicity so interwoven with each other that they
cannot be separated in an external manner. It was, after all, precisely
in the urge to gain greater clarity about their relation that we found
the possibility and the motive for the awakening of reflection.

At no element can we, in the domain of the mental life (any more than in
that of matter), come to a stop as absolutely isolated or absolutely
simple. There constantly poses itself, for reflection, the task of
finding a connection between the elements or objects that might seem to
be isolated, and of finding an inner multiplicity in the elements or
objects that appear as simple. It is a sign of the dullness of
reflection when it comes to a stop too early in either of these
respects. In other words: the law of synthesis must always be carried
through anew with respect to the objects or elements themselves.

On the other hand, neither can there be any absolutely conclusive
synthesis. As long as the mental life lasts, new objects are taken up
into the involuntary stream and are involuntarily worked into it. They
are new drops, or perhaps new brooks, which are partly reshaped by the
stream and partly can also gain influence on its character and its
course. To this is added the character of the new riverbeds that the
stream must form for itself when the terrain changes. And as long as
reflection is awake, it will constantly be able to set itself the task
of bringing about a greater, more comprehensive, and more intimate
synthesis than before. There is always a psychic work to be done. The
synthesis that has hitherto enclosed the multiplicity of the mental life
can again come to stand as object or element for a later synthesis. As
soon as we become conscious of our preceding mental states and
activities, such a transition to a new synthesis takes place, whether or
not this signifies a progress, a more valuable state. It is a new I that
becomes conscious of our earlier I's. No synthesis is definitive. Even
the philosophy that spoke most of higher unity always fell into
embarrassment when it had to speak of the highest unity.

It can of course be entirely justified, in a special connection, with a
view to a particular investigation, to disregard these difficulties of
principle and to regard elements or objects as if they were absolute and
simple. Every special scientific investigation presupposes a certain
abstraction. If, e.g., experimental psychology really needs to operate
with a kind of psychic atoms or flashes as independent and simple, this
can be just as justified as when, in natural science, one operates with
material atoms. But one must avoid the dogmatism that confuses concepts
we have arranged for the needs of an investigation with ultimate
realities. And likewise it can be justified that in everyday speech, or
in a practical---e.g. ethical or juridical---respect, we regard the
individual as an absolute unity with a completed synthesis, without
regard to the fact that synthesis at bottom signifies a work, an
activity, not something given once and for all. There is an irrational
relation between the mental life and what we call its elements. The
point now is to find as many elements as possible; but it can be
justified and necessary to confine oneself to a few decimal places,
provided only that one does not forget that the question is then not
exhausted. The irrational is here, as everywhere, both a barrier and a
task.

There expresses itself here an antinomy that is closely bound up with
the essence of the mental life and is peculiar to the concept of
personality. Mental life and personality cannot be conceived as products
of previously given elements, since the elements we know always have
their properties precisely by being members of the mental life. We here
come to move in a circle, unless we are clear that an inexhaustible
reciprocal action prevails: the parts are determined by the whole, and
the whole by the parts. In each particular case, in each particular
psychological consideration, we must investigate how far the synthesis,
the formation of the whole, has advanced, and on this basis observe how
the emergence of a new object operates. But the very distinction between
synthesis and object is always more or less artificial in relation to
the constant interplay of meeting and parting, of continuity and
discontinuity, of wholeness and dissolution. There is here a mental
rhythmics that forms an analogy to---or rather a prototype for---the
rhythms that characterize the course of organic life, a rhythmics that
by no means begins with reflection, but, by all we can judge, already
makes itself felt in the involuntary mental life. Indeed, this rhythmics
recurs in the very relation between reflection and the involuntary
course of the mental life. The difficulty of understanding the origin of
the mental life rests precisely on the fact that it cannot be explained
as a product of previously given elements, but from its very first
moment moves in a rhythm---and that in a rhythm of the peculiar kind
that none of the members in the rhythm can have absolute priority over
the other.

5. In the history of the concept of synthesis this antinomy has been a
stumbling-block. The concept of synthesis itself is an old idea,
traceable in Plato and Aristotle, and especially in Plotinus in
antiquity, and which from Leibniz's time on has always, with more or
less energy, been asserted in psychology and in philosophy
generally\footnote{On the history of the concept of synthesis in more
recent times, see \emph{Georges Dwelshauvers}: \emph{La synthèse
mentale} (Paris 1908), pp. 179--223 (an account of the concept of
synthesis in Leibniz, Kant, Wundt, Høffding, and Pierre Janet).---In
\emph{S. Kierkegaard} this concept appears in an energetic and profound
manner, especially in the measure he applies in his doctrine of the
stages. See my book \emph{Danske Filosofer}, pp. 153 f.; p. 165 f.}.

\emph{Hume} cut the knot by dissolving the mental life into absolutely
independent elements, and consistently he therefore found an insoluble
enigma in all unity and connection in consciousness. The principle of
unity (the uniting principle) was for him the unintelligible. \emph{Kant}
started from the opposite point of view, that it is precisely the
unconnected, the isolated and mutually independent, that is
unintelligible to us, both in psychology and wherever else we meet it.
But he still held fast to Locke's presupposition that, as the basis for
thought's combining and ordering activity, there was given a
multiplicity that was only now to be connected. Thereby he arrived at
his fateful distinction between the form of cognition and its matter.
This too was to cut the knot. And yet Kant did not lack the insight that
already in the involuntary mental life there operate forms analogous to
those that clear reflection employs. Synthesis is, he says, the effect of
a blind but indispensable function of the soul, without which we would
have no cognition at all, but of which we only rarely become conscious.
It is the task of thought to form the synthesis into a concept. And this
same blind function of the soul (Kant calls it imagination) is likewise
what conditions the connection between the two extremes of our cognition,
``sensation'' and ``understanding.'' It was surely only Kant's
reluctance to give psychological considerations too large a place in his
investigation that held him back from giving a fuller characterization
of the involuntary course of the mental life. J. F. \emph{Fries}
developed Kant's doctrine further from this side and is even regarded by
his disciples as the real discoverer of the involuntary
basis\footnote{Against this, see my review in \emph{Göttinger Gelehrte
Anzeiger} (1907) of ``Abhandlungen der Frieseschen Schule. Neue
Folge.''}. But he then believed that in this basis---the spontaneous
striving, the immediate cognition---he possessed a rule of truth exalted
above all error. Unfortunately this rule is ineffable, since it can come
to consciousness only in an imperfect and fragmentary way through
reflection! A dogmatism comes to the fore in the firm belief in the
validity of the involuntary mental life---a dogmatism that signifies a
decided step backward in relation to Kant. Already when Fries assumes
that ``spontaneity'' is not only continuous but also constant, he says
more than any reflection could demonstrate. The irrational relation
between immediacy and reflection is not strictly maintained.

This relation is, as already mentioned, very sharply emphasized by
\emph{Henri Bergson}. This thinker is not satisfied with the concept of
synthesis, because it always presupposes an opposition between unity and
multiplicity. Is it not an opposition---he asks---that recurs, whatever
one thinks of or speaks about? What matters is to say \emph{which} kind
of unity and \emph{which} kind of multiplicity we have in the mental
life. The streaming life of the personality cannot be defined by means
of concepts that always stand in a certain relation of opposition to one
another. One cannot, through thesis and antithesis, arrive at a real
unity. One must, conversely, immerse oneself in the immediate stream of
the mental life in order to see the opposition emerge from it. Immediate
beholding (intuition) must take the place of
reflection\footnote{\emph{Introduction à la Métaphysique}. (Revue de
Métaphysique et de Morale 1903, pp. 15--17.)}.

Bergson's standpoint recalls, viewed from one side, that of Fries,
except that he gives us, with far greater mastery, an understanding of
what is meant by ``spontaneity.'' But whereas Fries and his disciples
held that true cognition is attained by translating, as well as we can,
that spontaneous basis into the forms of reflection, Bergson demands that
we subdue reflection and, by an act of will, dive down again into the
immediate, feeling ourselves one with it. It stands here as an enigma
how the validity of a cognition can be demonstrated without reflection.
One of the special problems of thought will later give us occasion to
bring out this difficult point in the philosophy of the brilliant French
thinker.

In response to the remark that the concepts we use in our reflection to
characterize the mental life---synthesis, unity, multiplicity, and the
like---find application to all objects, I will maintain that the reason
why that group of concepts constantly finds, and must find, application
to all objects is to be sought precisely in the fact that they express
the forms in which the mental life, and cognition in particular, unfolds
itself. Only through the application of these forms does the
appropriation, the taking-up into the connection of our own being, in
which all understanding, all cognition consists, become possible. No
wonder that all products of cognition bear their stamp! Our cognition
consists in a conscious or unconscious analogy with our own being.---

The relation of opposition that has been discussed in the foregoing
under various expressions (object and reflection, matter and form,
spontaneity and reflection, intuition and analysis), \emph{Stumpf}
designates as a relation of opposition between phenomena and psychic
functions\footnote{\emph{Erscheinungen und psychische Funktionen}.
(Abhdl.\ der preussischen Gesellschaft der Wiss.\ 1907), pp. 15--21.}.
With respect to this relation he advances, with all possible
reservation, the hypothesis that psychic phenomena are not altered
because they are analyzed. When we become aware of something in an
observation that we did not notice at the first moment, this is supposed
to bring about no change in our observation itself, but only in the
degree of explicit consciousness we apply to it. A sound rings the same
whether or not I notice individual tones in it; a dish makes the same
impression whether I can distinguish something sour or sweet in its
taste, or whether the taste acts as a unified whole. In my judgment it
is, from the outset, not probable that this is how matters stand. The
differentiated and differentiating attention in which analysis consists
must alter the object for us. The fact that we have discovered a spot in
an otherwise beautiful picture---a spot we did not notice at the first
moment---can give the whole picture a different character. The timbre
changes; an element has, quite simply, been added that was not present
in the first glance. And the relation between the individual parts is in
each case a different one after the analysis than before. It can cost an
exertion of will to return to the immediate beholding in which the
particulars have not yet become the object of special apprehension.
There is here a difference between experience and reflection that can
scarcely be effaced. There are of course many degrees, as we shall come
to see in what follows, in the investigation of the relation between
intuiting and judging. And individual differences certainly also make
themselves felt. In elastic natures the transition from experience to
reflection, and vice versa, will take place relatively easily; and yet I
should think that even in them a certain contrast-effect will make
itself felt between the two states.

\emph{Stumpf}, for that matter, will only maintain the possibility of a
differentiated subsistence, during the experience, of what is first
differentiated with consciousness during reflection. And he concedes
that there is probably no absolute independence in the relation between
phenomenon and function ``under the usual, complicated circumstances of
the psychic process.'' Whether experimental investigations can, as he
surmises, lead to any result here is doubtful. It is not easy to see how
the difference between immediate experience and reflection could fall
away. And in a psychological experiment it is difficult to obtain a
wholly immediate experience.

Chemists often say that the atoms of the elements are present in
compound substances---so that, e.g., carbon and oxygen are really
present in carbonic acid. And Stumpf holds that the psychologist must
have the same right in his domain. I think, however, that there is a
difference. For the chemist can demonstrate that there are just as many
parts by weight in carbonic acid as in carbon and oxygen each by itself;
and at bottom that is all he means by that manner of speaking. Such an
equivalence the psychologist is, for good reasons, debarred from
demonstrating.---For that matter, a chemist such as \emph{Sir William
Ramsay} has expressly denied the right to take that manner of speaking
literally, and \emph{Ernst Mach} has expressed himself in the same
direction. The analogy is therefore doubtful.---

We do not get beyond an antinomy between experience and reflection, and
we may not, without further ado, transfer to the one what holds for the
other. Experience itself we may think of neither as absolute continuity
nor as a chaos of differences. The continuous and the discontinuous
occur in it in all possible degrees, and it is only reflection that
takes the opposition between continuity and discontinuity sharply into
view, in that it---as will later be shown---receives the double task of
finding differences where continuity was immediately given, and
intermediate links where discontinuity was immediately given.

% [translation continues from p. 20]

\section*{B. Intuition, Association, and Comparison}
\addcontentsline{toc}{section}{B. Intuition, Association, and Comparison}
\label{sec:anskuen}

% [text to be added: p. 37 — section intro, if any]

\subsection*{a. Sensation}
\addcontentsline{toc}{subsection}{a. Sensation}
\label{ssec:sansning}

% [text to be added: pp. 37--47]

\subsection*{b. Recognition}
\addcontentsline{toc}{subsection}{b. Recognition}
\label{ssec:genkendelse}

% [text to be added: pp. 47--52]

\subsection*{c. Memory- and Imagination-Intuition}
\addcontentsline{toc}{subsection}{c. Memory- and Imagination-Intuition}
\label{ssec:erindring}

% [text to be added: pp. 52--59]

\subsection*{d. Association of Ideas}
\addcontentsline{toc}{subsection}{d. Association of Ideas}
\label{ssec:association}

% [text to be added: pp. 59--64]

\subsection*{e. Comparison}
\addcontentsline{toc}{subsection}{e. Comparison}
\label{ssec:sammenligning}

% [text to be added: pp. 64--71]

\section*{C. Judging}
\addcontentsline{toc}{section}{C. Judging}
\label{sec:dommen}

% [text to be added: p. 71 — section intro, if any]

\subsection*{a. The Formation of Judgment}
\addcontentsline{toc}{subsection}{a. The Formation of Judgment}
\label{ssec:domsdannelse}

% [text to be added: pp. 71--74]

\subsection*{b. Obstacles to the Formation of Judgment}
\addcontentsline{toc}{subsection}{b. Obstacles to the Formation of Judgment}
\label{ssec:hindringer}

% [text to be added: pp. 74--82]

\subsection*{c. Subject and Predicate}
\addcontentsline{toc}{subsection}{c. Subject and Predicate}
\label{ssec:subjekt-praedikat}

% [text to be added: pp. 82--99]

\subsection*{d. The Validity of Judgment}
\addcontentsline{toc}{subsection}{d. The Validity of Judgment}
\label{ssec:gyldighed}

% [text to be added: pp. 99--109]

%% ============================================================
%%  II. THE HISTORY OF THOUGHT
%% ============================================================
\chapter*{II. The History of Thought}
\addcontentsline{toc}{chapter}{II. The History of Thought}
\label{ch:historie}

\section*{A. Animism, Platonism, and Positivism}
\addcontentsline{toc}{section}{A. Animism, Platonism, and Positivism}
\label{sec:animisme}

% [text to be added: pp. 109--123]

\section*{B. Ancient and Modern Thinking}
\addcontentsline{toc}{section}{B. Ancient and Modern Thinking}
\label{sec:antik-moderne}

% [text to be added: pp. 123--135]

%% ============================================================
%%  III. THE FORMS OF THOUGHT (THE CATEGORIES)
%% ============================================================
\chapter*{III. The Forms of Thought (The Categories)}
\addcontentsline{toc}{chapter}{III. The Forms of Thought (The Categories)}
\label{ch:former}

\section*{A. The History and Method of the Doctrine of Categories}
\addcontentsline{toc}{section}{A. The History and Method of the Doctrine of Categories}
\label{sec:kategorilaere-historie}

% [text to be added: pp. 135--154]

\section*{B. Fundamental Categories}
\addcontentsline{toc}{section}{B. Fundamental Categories}
\label{sec:fundamentale-kategorier}

% [text to be added: pp. 154--175]

\section*{C. Formal Categories}
\addcontentsline{toc}{section}{C. Formal Categories}
\label{sec:formale-kategorier}

% [text to be added: pp. 175--177 — section intro]

\subsection*{a. Identity}
\addcontentsline{toc}{subsection}{a. Identity}

% [text to be added: pp. 177--181]

\subsection*{b. Quality-Relations}
\addcontentsline{toc}{subsection}{b. Quality-Relations}

% [text to be added: pp. 181--187 — subsection intro]

\subsubsection*{1) Time}
\addcontentsline{toc}{subsubsection}{1) Time}

% [text to be added: pp. 187--190]

\subsubsection*{2) Number}
\addcontentsline{toc}{subsubsection}{2) Number}

% [text to be added: pp. 190--194]

\subsubsection*{3) Degree}
\addcontentsline{toc}{subsubsection}{3) Degree}

% [text to be added: pp. 194--197]

\subsubsection*{4) Place}
\addcontentsline{toc}{subsubsection}{4) Place}

% [text to be added: pp. 197--202]

\subsection*{c. Negation}
\addcontentsline{toc}{subsection}{c. Negation}

% [text to be added: pp. 202--204]

\subsection*{d. Rationality}
\addcontentsline{toc}{subsection}{d. Rationality}

% [text to be added: pp. 204--207]

\section*{D. Real Categories}
\addcontentsline{toc}{section}{D. Real Categories}
\label{sec:reale-kategorier}

% [text to be added: pp. 207--208 — section intro]

\subsection*{a. Causality}
\addcontentsline{toc}{subsection}{a. Causality}

% [text to be added: pp. 208--218]

\subsection*{b. Totality}
\addcontentsline{toc}{subsection}{b. Totality}

% [text to be added: pp. 218--227]

\subsection*{c. Development}
\addcontentsline{toc}{subsection}{c. Development}

% [text to be added: pp. 227--237]

\section*{E. Ideal Categories}
\addcontentsline{toc}{section}{E. Ideal Categories}
\label{sec:ideale-kategorier}

% [text to be added: pp. 237--242 — section intro]

\subsection*{1) Formal Differences}
\addcontentsline{toc}{subsection}{1) Formal Differences}

% [text to be added: pp. 242--244]

\subsection*{2) Real Differences}
\addcontentsline{toc}{subsection}{2) Real Differences}

% [text to be added: pp. 244--246]

%% ============================================================
%%  IV. THE TASKS OF THOUGHT (PROBLEMS)
%% ============================================================
\chapter*{IV. The Tasks of Thought (Problems)}
\addcontentsline{toc}{chapter}{IV. The Tasks of Thought (Problems)}
\label{ch:opgaver}

\section*{Introduction}
\addcontentsline{toc}{section}{Introduction}

% [text to be added: pp. 246--254]

\section*{A. Cognition}
\addcontentsline{toc}{section}{A. Cognition}
\label{sec:erkendelse}

% [text to be added: p. 254 — section intro, if any]

\subsection*{a. Forms and Principles}
\addcontentsline{toc}{subsection}{a. Forms and Principles}

% [text to be added: pp. 254--274]

\subsection*{b. Forms and Objects}
\addcontentsline{toc}{subsection}{b. Forms and Objects}

% [text to be added: pp. 274--275 — subsection intro]

\subsubsection*{α. Quality and Quantity}
\addcontentsline{toc}{subsubsection}{α. Quality and Quantity}

% [text to be added: pp. 275--278]

\subsubsection*{β. Succession and Causality}
\addcontentsline{toc}{subsubsection}{β. Succession and Causality}

% [text to be added: pp. 278--296]

\subsubsection*{γ. Subject and Object}
\addcontentsline{toc}{subsubsection}{γ. Subject and Object}

% [text to be added: pp. 296--306]

\section*{B. World-View}
\addcontentsline{toc}{section}{B. World-View}
\label{sec:verdensanskuelse}

% [text to be added: p. 306 — section intro, if any]

\subsection*{a. Science and World-View}
\addcontentsline{toc}{subsection}{a. Science and World-View}

% [text to be added: pp. 306--320]

\subsection*{b. Unity and Multiplicity}
\addcontentsline{toc}{subsection}{b. Unity and Multiplicity}

% [text to be added: pp. 320--330]

\subsection*{c. Spirit and Matter}
\addcontentsline{toc}{subsection}{c. Spirit and Matter}

% [text to be added: pp. 330--339]

\subsection*{d. Subsistence and Development}
\addcontentsline{toc}{subsection}{d. Subsistence and Development}

% [text to be added: pp. 339--347]

\section*{C. Valuation}
\addcontentsline{toc}{section}{C. Valuation}
\label{sec:vurdering}

% [text to be added: pp. 347--350 — section intro]

\subsection*{a. The Ethical Problem}
\addcontentsline{toc}{subsection}{a. The Ethical Problem}

% [text to be added: p. 350 — subsection intro, if any]

\subsubsection*{α. \emph{Ethical Work}. (Formal Ethics)}
\addcontentsline{toc}{subsubsection}{α. Ethical Work. (Formal Ethics)}

% [text to be added: pp. 350--359]

\subsubsection*{β. \emph{The Rationality of Ethical Valuation}. (Real Ethics)}
\addcontentsline{toc}{subsubsection}{β. The Rationality of Ethical Valuation. (Real Ethics)}

% [text to be added: pp. 359--369]

\subsection*{b. The Religious Problem}
\addcontentsline{toc}{subsection}{b. The Religious Problem}

% [text to be added: pp. 369--374 — subsection intro]

\subsubsection*{α. \emph{The Psychological Place of Religion}}
\addcontentsline{toc}{subsubsection}{α. The Psychological Place of Religion}

% [text to be added: pp. 374--385]

\subsubsection*{β. \emph{The Historical and the Philosophical Consideration}}
\addcontentsline{toc}{subsubsection}{β. The Historical and the Philosophical Consideration}

% [text to be added: pp. 385--392]

\end{document}
